\begin{lemma}[The invertible binary pencil case]
\label{B:binary-pencil-Tor}
Let $k$ be either $\mathbb Q$ or $\mathbb C$, and let $A_B$ be the
four-dimensional $k$-algebra with basis $1,x,y,c$ and multiplication
\[
 x_i x_j=b_{ij}c,\qquad (x_1,x_2)=(x,y),\qquad
 c^2=cx=xc=cy=yc=0,
\]
where $B=(b_{ij})\in\operatorname{Mat}_2(\{0,1\})$ is invertible.
For finite-dimensional right and left $A_B$-modules $W,V$, respectively,
if $\operatorname{Tor}_2^{A_B}(W,V)\ne0$, then
$\operatorname{Tor}_n^{A_B}(W,V)\ne0$ in arbitrarily high degrees.
\end{lemma}

\begin{proof}
We first work over $k=\mathbb C$ and write $A=A_B$ and
$J=\langle x,y,c\rangle$.  Then $J^2=kc$, $J^3=0$, and $A/J=k$.
The bilinear pairing
\[
 A\times A\longrightarrow k,\qquad
 (a,b)\longmapsto\text{the coefficient of $c$ in $ab$}
\]
is nondegenerate: it pairs $1$ with $c$ in both directions, and its
restriction to $\langle x,y\rangle$ has matrix $B$.
Consequently $A$ is Frobenius and is injective as both a left and a
right module over itself.  More explicitly, the pairing identifies
the appropriate regular module with $\operatorname{Hom}_k(A,k)$,
which is injective because its module Hom functor is the exact
vector-space dual functor.  The left and right socles of $A$ are
both $kc$.  In particular, every nonzero left ideal contains $kc$.

We describe the nonprojective indecomposable left modules.
If $M$ is such a module and $cm\ne0$ for some $m\in M$, then
the homomorphism $A\to M$, $a\mapsto am$, is injective: any
nonzero kernel would contain $c$.  Its image splits because $A$
is injective, contradicting the assumed indecomposability and
nonprojectivity of $M$.  Thus
\[
 cM=0.                                                     \tag{\ref{B:binary-pencil-Tor}.1}\label{eq:B-binary-c-kills}
\]
Moreover, unless $M=k$, one has $\operatorname{soc}M=JM$.
Indeed, $JM\subseteq\operatorname{soc}M$ by
\eqref{eq:B-binary-c-kills}.  A socle
vector outside $JM$ spans a simple direct summand: choose a
vector-space complement containing $JM$, which is automatically
an $A$-submodule since $J$ sends the whole module into $JM$.

Choose a complement $T$ to $S=JM$ in $M$.  The actions of $x,y$
are therefore a pair of linear maps
\[
 X,Y:T\longrightarrow S,
\]
with $\operatorname{im}X+\operatorname{im}Y=S$ and, for $M\ne k$,
$\ker X\cap\ker Y=0$.
Isomorphism of modules is exactly simultaneous equivalence of
these two matrices, and a direct-sum decomposition of the pencil
is exactly a direct-sum decomposition of the module.
The Kronecker canonical form gives the following complete list
of indecomposable possibilities over $\mathbb C$:
\[
\begin{array}{c|c|c}
 \text{module}&(\dim T,\dim S)&(X,Y)\\ \hline
 I_n&(n+1,n)&([I_n\ 0],[0\ I_n]),\quad n\geq1\\
 P_n&(n,n+1)&([I_n\ 0]^{\mathsf T},[0\ I_n]^{\mathsf T}),
                         \quad n\geq1\\
 R_{n,\lambda}&(n,n)&(I_n,J_n(\lambda)),
                         \quad n\geq1,\ \lambda\in\mathbb C\\
 R_{n,\infty}&(n,n)&(J_n(0),I_n),\quad n\geq1.
\end{array}                                                \tag{\ref{B:binary-pencil-Tor}.2}\label{eq:B-binary-pencil-list}
\]
There is also the simple module $k$, with dimensions $(1,0)$.
The other zero-size singular pencil, of dimensions $(0,1)$,
is excluded by $S=JM$.
Here $J_n(\lambda)$ is a single Jordan block.  The exact matrix
pencil theorem used in this list is recorded, with these singular
and regular blocks, in
\cite[Theorem~2.1, p.~221]{deterandmytryshyndopico2024}.

We next determine syzygies, keeping track of indecomposability.
If $a=\dim(M/JM)$, the minimal projective cover is $A^a\to M$.
Its kernel $\Omega M$ lies in $JA^a$, is killed by $c$, and
contains $cA^a=\operatorname{soc}(A^a)$.  Thus $A^a$ is the
injective envelope of $\Omega M$.  If $\Omega M$ were a direct
sum of two nonzero modules, its injective envelope, and hence
the cokernel $M$, would split accordingly.  Neither resulting
cokernel can be zero: a nonzero injective summand is a sum of
copies of $A$, whereas $c$ kills $\Omega M$.
It follows that $\Omega M$ is again indecomposable and
nonprojective.

Put $b=\dim JM$.  The socle of $\Omega M$ has dimension $a$
and $\dim\Omega M=3a-b$.  Unless $\Omega M=k$, its top and
radical dimensions are consequently $(2a-b,a)$.
If $\Omega M=k$, these same dimension equalities force
$(a,b)=(1,2)$, so $M=P_1$.
The list \eqref{eq:B-binary-pencil-list}, and the direct
calculation $\Omega k=J\cong I_1$,
now give
\[
 \Omega k\cong I_1,\qquad
 \Omega I_n\cong I_{n+1},\qquad
 \Omega P_1\cong k,\qquad
 \Omega P_n\cong P_{n-1}\quad(n\geq2).                     \tag{\ref{B:binary-pencil-Tor}.3}\label{eq:B-binary-string-syzygies}
\]
Thus all the singular-pencil modules are positive or negative
syzygies of the simple module; the displayed formulas suffice
without choosing a complete resolution.

For the regular pencils there is an explicit formula that also
proves periodicity.  Identify $T=S=k^n$, and initially take
$(X,Y)=(I_n,J_n(\lambda))$ or $(J_n(0),I_n)$.
More generally the following calculation applies whenever $X,Y$
commute and at least one is invertible.  In the cover
$A\otimes_k T\to M$, the kernel is the direct vector-space sum
of $c\otimes T$ and the vectors
\[
 q(v)=y\otimes Xv-x\otimes Yv,\qquad v\in T.
\]
Indeed, modulo $c\otimes T$ its kernel is that of
$(u,v)\mapsto Xu+Yv$, and the displayed vectors form a basis
of that kernel.  Using $q(v)$ as the top coordinates and
$c\otimes v$ as the radical coordinates, multiplication gives
\[
 \Omega(X,Y)=
 \bigl(b_{12}X-b_{11}Y,\ b_{22}X-b_{21}Y\bigr).             \tag{\ref{B:binary-pencil-Tor}.4}\label{eq:B-binary-regular-syzygy}
\]
In particular, if
\[
 D_B=\begin{pmatrix}b_{12}&-b_{11}\\ b_{22}&-b_{21}\end{pmatrix},
\]
the $r$th syzygy is represented by $D_B^r(X,Y)^{\mathsf T}$.
This iteration is legitimate in the same coordinates.  Both
matrices remain commuting linear combinations of the original
$I_n$ and its single Jordan block.  Since $D_B$ is invertible,
their eigenvalues cannot both vanish, so at least one of the
two matrices remains invertible at every step.

There are six invertible binary $2$ by $2$ matrices.
Four are symmetric; for each of them
$\operatorname{tr}D_B=0$ and $\det D_B=\det B$, whence
\[
 D_B^2=-\det(B)I_2.
\]
The other two, and their powers, are
\[
\begin{split}
 B=\begin{pmatrix}1&1\\0&1\end{pmatrix}:&\quad
 D_B=\begin{pmatrix}1&-1\\1&0\end{pmatrix},\quad D_B^3=-I_2,\\
 B=\begin{pmatrix}1&0\\1&1\end{pmatrix}:&\quad
 D_B=\begin{pmatrix}0&-1\\1&-1\end{pmatrix},\quad D_B^3= I_2.
\end{split}                                               \tag{\ref{B:binary-pencil-Tor}.5}\label{eq:B-binary-periods}
\]
Multiplying both matrices of a pencil by the same nonzero
scalar does not change its module isomorphism type, since one
may rescale the radical coordinates.  Every regular
indecomposable is therefore periodic under $\Omega$, with
period dividing two in the symmetric case and dividing three
in the other two cases.  The right-module calculation is the
same calculation for $A^{\mathrm{op}}=A_{B^{\mathsf T}}$.

Finally decompose $W,V$ into indecomposable summands and select
a pair with nonzero degree-two Tor.  Neither summand is
projective.  If either is a regular-pencil module of period
$q$, dimension shifting gives
\[
 \operatorname{Tor}_{2+rq}^{A}(W,V)
       \cong\operatorname{Tor}_{2}^{A}(W,V)\ne0
       \qquad(r\geq0)
\]
for the selected pair.  Suppose instead that both are
singular-pencil modules or simple modules.  By
\eqref{eq:B-binary-string-syzygies} there exist
nonnegative integers $r_W,s_W,r_V,s_V$ with
\[
 \Omega^{r_W}W\cong\Omega^{s_W}k,
 \qquad
 \Omega^{r_V}V\cong\Omega^{s_V}k.
\]
For all $n>r_W+r_V$, ordinary dimension shifting on the two
sides yields
\[
 \operatorname{Tor}_n^A(W,V)
 \cong
 \operatorname{Tor}_{n-r_W-r_V+s_W+s_V}^A(k,k).
\]
Every positive-degree group on the right is nonzero.  For
example, \eqref{eq:B-binary-string-syzygies} shows directly that the minimal free resolution
of $k$ never terminates, and all its differentials become zero
after tensoring with $k$.  This proves the required recurrence
over $\mathbb C$.

For $k=\mathbb Q$, extension of scalars along
$\mathbb Q\subset\mathbb C$ gives, in every degree,
\[
 \mathbb C\otimes_{\mathbb Q}\operatorname{Tor}_n^{A_B}(W,V)
 \cong
 \operatorname{Tor}_n^{\mathbb C\otimes A_B}
       (\mathbb C\otimes W,\mathbb C\otimes V).
\]
This follows by extending any free resolution, and the extension
is faithfully flat.  Nonzero degree-two Tor therefore survives
over $\mathbb C$, where the result just proved supplies nonzero
groups in arbitrarily high degrees.  Faithfulness brings their
nonvanishing back to $\mathbb Q$.
\end{proof}
